\documentclass[10pt,a4paper]{amsart}
\usepackage[margin=19mm]{geometry}
\usepackage{amsmath,amssymb,mathtools}
\usepackage[scr=rsfs,cal=euler]{mathalfa}
\usepackage{xfrac}
\usepackage[unicode,hidelinks,bookmarks=false]{hyperref}
\newcommand{\ie}{i.e.\ }
\newcommand{\cf}{cf.\ }
\newcommand{\resp}{respectively\ }
\newcommand{\Subsection}{\subsection}
\providecommand{\nobreakdash}{\nobreak\hbox{-}}
\setlength{\emergencystretch}{2em}
\multlinegap0pt
\usepackage{amssymb}
\newtheorem{theorem}{Theorem}
\title{\bf Taut and Dupin Submanifolds (Updated Version)}
\author{Thomas E. Cecil}
\date{Source: arXiv:2609.12152v1 (CC BY 4.0)}
\begin{document}
\maketitle

\noindent\emph{Source and licence.} \bf Taut and Dupin Submanifolds (Updated Version). Version arXiv:2609.12152v1 (CC BY 4.0), distributed by arXiv under the Creative Commons Attribution 4.0 International licence, http://creativecommons.org/licenses/by/4.0/, as recorded in the arXiv licence metadata for this exact version and shown on the abstract page https://arxiv.org/abs/2609.12152v1. Full licence notice: \texttt{licenses/arxiv-2609.12152-CC-BY-4.0.txt}.\par\smallskip

\noindent\textit{Editorial scope.} Counted unit: the article's theorem thm:4.1.1 (source0.tex lines 1245--1252). The article's complete original proof of it is delivered with it (source0.tex lines 1253--1279, one \texttt{proof} environment of 27 lines). No mathematics is written, repaired or re-derived: the statement and the proof are the article's own lines, in the article's own notation, and no retained line is substituted or re-flowed.  No further source-local context is needed: every cross-reference of every retained line resolves inside the two delivered spans.  Nothing was omitted from any retained span, so the package carries no omission record.  No retained line contains a citation, so the wrapper carries no bibliography and no bibliography entry.  No retained line contains a figure environment, an image-inclusion command or drawing code, so no image is delivered and the package declares no asset dependency.  Editorial formatting only, in addition: an AMS \texttt{amsart} preamble that restates the article's own declarations and macros used by the retained lines, namely the environments \texttt{theorem} on separate counters, together with the standard packages those lines need.  The retained environments print in this document as Theorem 1 in order of appearance, whereas the article prints the same items with the numbering of its own full document.  The complete native source of the article as distributed in the arXiv e-print for this exact version is bundled unchanged as \texttt{sources/210085/source0.tex}.
\begin{theorem}
\label{thm:4.1.1}
Given positive integers $m_1, \ldots , m_g$, with 
$m_1 + \cdots + m_g = n-1$, there exists a proper
Dupin hypersurface $M^{n-1}$ in ${\bf R}^n$ with $g$
distinct principal curvatures having respective
multiplicities $m_1, \ldots , m_g$.
\end{theorem} 

\begin{proof}
The proof is by an inductive construction
which will be clear once the first few cases are handled.
First we construct a proper Dupin hypersurface $M^3$ in 
${\bf R}^4$ with three principal curvatures of multiplicity
one. Begin with an open subset $U$ of a torus of
revolution in ${\bf R}^3$ on which neither principal
curvature vanishes.  Take $M^3$ to be the cylinder
$U \times {\bf R}$ in ${\bf R}^3 \times {\bf R} =
{\bf R}^4$.  Then $M^3$ has three distinct principal
curvatures at each point, one of which is identically
zero.  These are clearly constant along their
corresponding lines of curvature.
Next, to construct a proper Dupin hypersurface in ${\bf R}^5$
with three distinct principal curvatures having
respective multiplicities 1, 1, 2, take a cylinder
$U \times {\bf R}^2$ in ${\bf R}^3 \times {\bf R}^2$. Finally,
to get a proper Dupin hypersurface $V^4$ in
${\bf R}^5$ with four principal curvatures of multiplicity 1, first 
invert the hypersurface $M^3$ constructed above in a 3-sphere
in ${\bf R}^4$ chosen so that the image of $M^3$
contains an open set $W^3$ on which no principal
curvature vanishes.  The hypersurface $W^3$ is proper Dupin, since the proper Dupin property is preserved by
M\"{o}bius transformations.
Now take $V^4$ to be the
cylinder $W^3 \times {\bf R}$. 
\end{proof}

\end{document}
